\subsection{Positive elements in stellar algebras}

DEFINITION 6. --- Let A be a stellar algebra. An element x of A is said to be positive if it is Hermitian and if \(\mathrm{Sp}_{\mathrm{A}}^{\prime}(x) \subset \mathbf{R}_{+}\) . One denotes by \(\mathrm{A}_{+}\) the set of positive elements of A. It is a subset of \(\mathrm{A}_{h}\) .

We denote \(x \geqslant y\) if \(x - y \in \mathrm{A}_{+}\) .

If the stellar algebra A is unital, its identity element is positive.

If B is a stellar subalgebra of A, one has \(B_{+} = B \cap A_{+}\) (cor. of prop. 5 of I, p.~106).

If \(\pi\colon A\to B\) is a morphism of stellar algebras, then \(\pi(A_{+})\subset B_{+}\) .

Examples. --- 1) Let X be a locally compact space. In the stellar algebra \(\mathcal{C}_{0}(\mathrm{X})\) of continuous functions on X tending to 0 at infinity, resp. in the stellar algebra \(\mathcal{C}_{b}(\mathrm{X})\) of bounded continuous functions on X, a function \(f\) is a positive element if and only if it is real-valued and if \(f(x) \geqslant 0\) for all \(x \in \mathrm{X}\) (cf.~example 3 of I, p.~17).

\begin{enumerate}
\def\labelenumi{\arabic{enumi})}
\setcounter{enumi}{1}
\item
  Let A be a commutative stellar algebra. Let \(a\) be an element of A. Since \(\mathrm{Sp}_{\mathrm{A}}^{\prime}(x)\) is the union of \(\{0\}\) and of the image of the Gelfand transform \(\mathcal{G}(a)\) (prop. 6 of I, p.~37), the element \(a\) is positive if, and only if, the Gelfand transform \(\mathcal{G}(a)\) is a positive function.
\item
  Let E be a complex Hilbert space. An element x of the stellar algebra \(\mathcal{L}(\mathrm{E})\) (example 1 of I, p.~102) is positive if and only if it is a positive endomorphism of E in the sense of EVT, V, p.~45, def. 6 (prop. 8 of I, p.~138).
\end{enumerate}

Lemma 12. --- Let A be a unital stellar algebra and let \(x \in A\) be a Hermitian element.

\begin{enumerate}
\def\labelenumi{\alph{enumi})}
\item
  The element \(x\) is positive if and only if \(\left\| \|x\| \cdot 1 - x \right\| \leqslant \|x\|\) ;
\item
  If \(\| x\| \leqslant 1\) , then \(x\) is positive if and only if \(\| 1 - x\| \leqslant 1\) .
\item
  If x is positive, then 1 - x is positive if and only if \(\|x\| \leqslant 1\) ;
\item
  If x is positive and if \(y \in A_{+}\) commutes with x, then xy is positive.
\end{enumerate}

The element \(x\) is Hermitian, hence normal. By considering the stellar subalgebra generated by \(x\) , which is commutative, one reduces to the case where the algebra A is commutative, that is, to the case where \(A = \mathcal{C}_0(X)\) for a locally compact topological space X (th. 1 of I, p.~108). The first three assertions then follow immediately from example 1 above. Likewise, to prove assertion \(d\) ), one may consider the stellar subalgebra generated by \(x\) and \(y\) , which is commutative.

PROPOSITION 14. --- Let A be a stellar algebra. The set \(\mathrm{A}_{+}\) is a closed pointed salient convex cone in the real Banach space \(\mathrm{A}_{h}\) (EVT, II, p.~11).

Let \(\widetilde{\mathsf{A}}\) be the stellar algebra deduced from \(\mathsf{A}\) by adjoining an identity element. Since \(\mathsf{A}_{+} = \mathsf{A} \cap \widetilde{\mathsf{A}}_{+}\) , it suffices to prove the proposition for \(\widetilde{\mathsf{A}}\) . One may therefore suppose that \(\mathsf{A}\) has a unit element.

We have \(0 \in \mathrm{A}_{+}\) . For every \(\lambda \in \mathbf{R}_{+}^{*}\) and every \(x \in \mathrm{A}\) , one has \(\mathrm{Sp}_{\mathrm{A}}^{\prime}(\lambda x) = \lambda \mathrm{Sp}_{\mathrm{A}}^{\prime}(x)\) which implies that \(\mathrm{A}_{+}\) is a cone in the real Banach space \(\mathrm{A}_{h}\) .

To show that \(A_{+}\) is convex, it suffices to show that if x and y are positive, then \(x + y \geqslant 0\) (EVT, II, p.~11, prop. 10). By homothety, it suffices to prove that if \(x \geqslant 0\) and \(y \geqslant 0\) additionally satisfy \(\|x\| \leqslant 1\) , \(\|y\| \leqslant 1\) , then the element \(\frac{1}{2}(x + y)\) is positive. But we have

\[
\left\| 1 - \frac {1}{2} (x + y) \right\| \leqslant \frac {1}{2} \| 1 - x \| + \frac {1}{2} \| 1 - y \| \leqslant 1
\]

according to the assertion \(b)\) of Lemma 12, and this same assertion then shows that \(\frac{1}{2}(x + y)\) is positive.

Finally, the assertion \(a\) ) of Lemma 12 also implies that \(\mathrm{A}_{+}\) is closed.

Since \(A_{+}\) is a pointed cone in \(A_{h}\) , it is salient if and only if \(\mathrm{A}_{+}\cap(-\mathrm{A}_{+})\) is reduced to 0. But if \(x\in\mathrm{A}_{+}\cap(-\mathrm{A}_{+})\) , one has \(\mathrm{Sp}_{\mathrm{A}}^{\prime}(x)=\{0\}\) , hence \(\varrho(x)=0\) , and \(\|x\|=0\) . Since x is Hermitian (Lemma 7, (2) of I, p.~104), it follows that x=0.

Proposition 14 means that the relation \(x\geqslant y\) is an order relation on \(\mathrm{A}_h\) (EVT, II, p.~13, prop. 13).

PROPOSITION 15. — Let A be a stellar algebra. Let x be a normal element of A.

\begin{enumerate}
\def\labelenumi{\alph{enumi})}
\item
  Suppose that A is unital and let f be a continuous function from \(\mathrm{Sp}_{\mathrm{A}}(x)\) to C. For \(f(x)\) to be positive, it is necessary and sufficient that the image of f be contained in \(R_{+}\) ;
\item
  Let f be a continuous function from \(\mathrm{Sp}_{\mathrm{A}}^{\prime}(x)\) to C such that \(f(0)=0\) . For \(f(x)\) to be a positive element of A, it is necessary and sufficient that the image of f be contained in \(R_{+}\) .
\end{enumerate}

Assertion \(a\) ) follows from assertion \(a\) ) of Cor. 2 of I, p.~111, and assertion \(b\) ) follows from Prop. 13 of I, p.~114.

Let \(x\) be a Hermitian element of the stellar algebra A. Its spectrum is contained in \(\mathbf{R}\) (Prop. 4 of I, p.~106). Consider the continuous functions from \(\mathrm{Sp}_{\mathrm{A}}^{\prime}(x)\) into \(\mathbf{R}\) defined by

\[
f _ {1} \colon t \mapsto \sup (t, 0), \quad f _ {2} \colon t \mapsto \sup (- t, 0), \quad f _ {3} \colon t \mapsto | t |.
\]

We denote

\[
x ^ {+} = f _ {1} (x), \quad x ^ {-} = f _ {2} (x), \quad | x | = f _ {3} (x).\tag{3}
\]

Since the functions \(f_1, f_2, f_3\) are real-valued, positive, and vanish at 0, the elements \(x^+, x^-\) and \(|x|\) are positive elements of A (Prop. 15, a)) which belong to the stellar subalgebra of A generated by \(x\) .

We have \(f_{1}(t) - f_{2}(t) = t\) for all \(t \in \mathbf{R}\) , as well as the relations \(f_{1} + f_{2} = f_{3}\) and \(f_{1}f_{2} = 0\) . It follows that

\[
x = x ^ {+} - x ^ {-}, \quad | x | = x ^ {+} + x ^ {-}, \quad x ^ {+} x ^ {-} = x ^ {-} x ^ {+} = 0.\tag{4}
\]

Since the functional calculus map is isometric, we have

\[
\| | x | \| = \| x \|, \quad \| x ^ {+} \| \leqslant \| x \|, \quad \| x ^ {-} \| \leqslant \| x \|.
\]

Let \(x\) be a positive element of A. It is Hermitian, hence normal. Let \(\alpha \in \mathbf{R}_{+}^{*}\) , and let \(g\) be the restriction to \(\mathrm{Sp}_{\mathrm{A}}^{\prime}(x)\) of the function \(t \mapsto t^{\alpha}\) ; we denote \(x^{\alpha} = g(x)\) . This is a positive element of the stellar subalgebra of A generated by \(x\) . Let \(\alpha\) and \(\beta\) be in \(\mathbf{R}_{+}^{*}\) . Since the functional calculus map is an algebra homomorphism, and by Prop. 13 of I, p.~114, we have

\[
x ^ {\alpha} x ^ {\beta} = x ^ {\alpha + \beta} \quad (x ^ {\alpha}) ^ {\beta} = x ^ {\alpha \beta}.\tag{5}
\]

We shall also write \(\sqrt{x} = x^{1/2}\) .

PROPOSITION 16. — Let x be a positive element of A. Let \(\alpha \in \mathbf{R}_{+}^{*}\) . Then \(x^{1/\alpha}\) is the unique positive element y of A such that \(y^{\alpha} = x\) .

We saw above that \(x^{1/\alpha}\) satisfies the required properties. Conversely, let y be a positive element of A such that \(y^{\alpha} = x\) . By formula (5), we have \(y = (y^{\alpha})^{1/\alpha} = x^{1/\alpha}\) , which is what had to be proved.

Lemma 13. — Let A be a unital stellar algebra. Every element of A is a sum of unitary elements.

Let \(x\) be a Hermitian element of A. Suppose first that \(\| x \| \leqslant 2\) . By Lemma 12, c), we have \(1 - \frac{1}{4} x^2 \in \mathrm{A}_+\) . Let \(y = \frac{1}{2} x + i\sqrt{1 - \frac{1}{4} x^2}\) . We have \(y^* = \frac{1}{2} x - i\sqrt{1 - \frac{1}{4} x^2}\) , hence \(yy^* = 1\) and \(x = y + y^*\) is a sum of two unitary elements. In the general case, let \(k\) be an integer such that \(\| \frac{1}{k} x \| \leqslant 2\) ; the element \(x\) is then a sum of \(2k\) unitary elements. By Lemma 2 of I, p.~96, the lemma follows.

THEOREM 2. — Let A be a stellar algebra. An element x of A is positive if and only if there exists \(y \in A\) such that \(x = y^{*}y\) .

Suppose that \(x\) is positive. Let \(y = x^{1/2}\) ; this is a Hermitian element of A and we have \(y^*y = y^2 = x\) .

Conversely, let y be an element of A and set \(x = y^{*}y\) . This is a Hermitian element of A. Let us show that x is positive. For this, denote \(z = x^{-}\) and set w = yz. We then have

\[
w ^ {*} w = z ^ {*} y ^ {*} y z = z x z = z (x ^ {+} - z) z = - z ^ {3}.
\]

Since \(z \geqslant 0\) , one has \(z^3 \geqslant 0\) and we deduce that \(\mathrm{Sp}_{\mathrm{A}}'(w^{*}w) \subset \mathbf{R}_{-}\) . Let us also write \(w = w_1 + iw_2\) where \(w_1\) and \(w_2\) are Hermitian (lemma 2 of I, p.~96). The elements \(w_1^2\) and \(w_2^2\) are positive. We have \(ww^{*} + w^{*}w = 2w_{1}^{2} + 2w_{2}^{2}\) , and prop. 14 therefore shows that \(ww^{*} = 2w_{1}^{2} + 2w_{2}^{2} + (-w^{*}w)\) is positive. Since \(\mathrm{Sp}_{\mathrm{A}}'(ww^{*}) = \mathrm{Sp}_{\mathrm{A}}'(w^{*}w)\) (I, p.~5, prop. 1), which is contained in \(\mathbf{R}_{-}\) , we conclude that \(\mathrm{Sp}_{\mathrm{A}}'(w^{*}w) = \{0\}\) . Since \(w^{*}w\) is Hermitian, this implies (cor. 1 of I, p.~108) that \(\|w^{*}w\| = \varrho(w^{*}w) = 0\) , so that \(z^{3} = 0\) . Since z is Hermitian, we have z = 0. Thus, \(x = x^{+}\) is positive.

Remark. --- Let A be a stellar algebra and let \(x \in A\) . The element \(x^{*}x\) is positive; we therefore set \(|x| = (x^{*}x)^{1/2}\) . We have \(\|x\|^{2} = \|x^{*}x\| = |||x|^{2}\| = |||x||^{2}\) , hence \(\|x\| = |||x|\|\) .

When x is normal, we also have \(|x| = f(x)\) , where f is the map from C to C, vanishing at 0, given by \(f(z) = |z|\) . In particular, when x is Hermitian, \(|x|\) coincides with the element defined by formula (3).

Suppose further that A is unital and that x is invertible. Then \(|x|\) is also invertible, the element \(u = x|x|^{-1}\) is unitary and we have \(x = u|x|\) (polar decomposition; see also I, p.~139, no. 8 for the case of endomorphisms of Hilbert spaces).

Lemma 14. --- Let A be a stellar algebra, and let x and y be Hermitian elements of A.

\begin{enumerate}
\def\labelenumi{\alph{enumi})}
\item
  If \(x \leqslant y\) , then for every element \(w\) of A, we have \(w^{*}xw \leqslant w^{*}yw\) . In particular, if \(y \geqslant 0\) , one has \(w^{*}yw \geqslant 0\) ;
\item
  Suppose that A is unital. If \(0 \leqslant x \leqslant y\) and if \(x\) is invertible, then \(y\) is invertible and \(y^{-1} \leqslant x^{-1}\) ;
\item
  If \(0 \leqslant x \leqslant y\) then \(\|x\| \leqslant \|y\|\) .
\end{enumerate}

Let \(u = (y - x)^{1/2}\) . We have \(w^{*}yw - w^{*}xw = w^{*}u^{2}w = (uw)^{*}(uw)\) and the assertion a) follows from th. 2.

Let us prove assertion \(b\) ). Suppose first that \(x = 1\) . Let B be the unital stellar subalgebra generated by \(y\) . By the Gelfand isomorphism, \(y\) corresponds to a continuous function \(\geqslant 1\) on the compact space \(\mathsf{X}(\mathsf{B})\) . This function is therefore invertible, and its inverse is \(\leqslant 1\) . This implies that \(y\) is invertible and \(y^{-1} \leqslant 1 = x^{-1}\) . In the general case, we observe that \(0 \leqslant 1 \leqslant x^{-1/2}yx^{-1/2}\) according to \(a\) ), so the preceding case implies that \(z = x^{-1/2}yx^{-1/2}\) is invertible and \(z^{-1} \leqslant 1\) . Thus, \(y\) is invertible and \(y^{-1} \leqslant x^{-1}\) according to \(a\) ) again.

To prove the assertion \(c\) ), we may assume that A is unital (prop. 3 of I, p.~105). Suppose first that \(y\) is invertible. Let \(b = \sqrt{y}\) . According to \(a\) ), the conditions \(0 \leqslant x \leqslant y\) imply \(0 \leqslant b^{-1}xb^{-1} \leqslant b^{-1}yb^{-1} = 1\) , hence \(\|b^{-1}xb^{-1}\| \leqslant 1\) by Lemma 12, \(c\) ) of I, p.~116. We then have

\[
\| x \| = \| b (b ^ {- 1} x b ^ {- 1}) b \| \leqslant \| b \| \| b ^ {- 1} x b ^ {- 1} \| \| b \| \leqslant \| b \| ^ {2} = \| b ^ {2} \| = \| y \|.
\]

In the general case, for every real \(\varepsilon > 0\) , the element \(y + \varepsilon\) is invertible and \(0 \leqslant x \leqslant y + \varepsilon\) . From what precedes, we therefore have \(\|x\| \leqslant \|y + \varepsilon\|\) for every real number \(\varepsilon > 0\) , whence the result.

Note. --- In general, if x and y are positive elements of a stellar algebra A, the condition \(0 \leqslant x \leqslant y\) does not imply \(x^{2} \leqslant y^{2}\) (cf. exercise 15 of I, p. 182).

